\documentclass[10pt,a4paper]{amsart}
\usepackage[margin=19mm]{geometry}
\usepackage{amsmath,amssymb,mathtools}
\usepackage[scr=rsfs,cal=euler]{mathalfa}
\usepackage{xfrac}
\usepackage[unicode,hidelinks,bookmarks=false]{hyperref}
\newcommand{\ie}{i.e.\ }
\newcommand{\cf}{cf.\ }
\newcommand{\resp}{respectively\ }
\newcommand{\Subsection}{\subsection}
\providecommand{\nobreakdash}{\nobreak\hbox{-}}
\setlength{\emergencystretch}{2em}
\multlinegap0pt
\usepackage{mathtools}
\usepackage{enumitem}
\usepackage{color}
\newtheorem{theorem}{Theorem}[section]
\newtheorem{proposition}[theorem]{Proposition}
\newtheorem{lemma}[theorem]{Lemma}
\newtheorem{corollary}[theorem]{Corollary}
\newtheorem{remark}[theorem]{Remark}
\newtheorem{example}[theorem]{Example}
\newcommand{\Hessp}{\operatorname{Hess}^{+}}
\newcommand{\R}{\mathbb{R}}
\newcommand{\ip}[2]{\left\langle #1,#2\right\rangle}
\newcommand{\norm}[1]{\left\lVert #1\right\rVert}
\begin{document}
\title[High-level convexity for products of squared Euclidean dista]{High-level convexity for products of squared Euclidean distance functions}
\author{Tudor Micu, Cornel Pintea, George C. Ţurcaş}
\date{}
\maketitle
\noindent\emph{Source and licence.} High-level convexity for products of squared Euclidean distance functions. Version arXiv:2606.00316v1 (CC BY 4.0). This material is distributed under the Creative Commons Attribution 4.0 International licence, http://creativecommons.org/licenses/by/4.0/, as recorded in the arXiv OAI licence metadata for this exact version and shown on the abstract page https://arxiv.org/abs/2606.00316v1. Full rights notice: licenses/arxiv-2606.00316-CC-BY-4.0.txt.\par\smallskip
\noindent\textit{Editorial scope.} Counted unit: the original \texttt{theorem} environment \texttt{thm:two-centre-hess-region} at source lines 479--486 together with its complete proof at lines 488--520; the delivered document keeps source lines 404--521 so that every referenced label and every citation resolves inside it. Native editable source is the paper's own concatenated TeX; the AMS wrapper is editorial formatting only. No publisher class, upstream script or unrelated artwork is redistributed.
\par\smallskip\noindent\textit{Reference note.} This article ships no compiled bibliography (biblatex); the entries delivered here are composed from the author's own BibTeX (.bib) fields, with no field invented and no entry dropped.
\section{The two-centre product}

Let $p\ne q$ and consider the function
\[
 F_{p,q}(x)=\norm{x-p}^2\norm{x-q}^2.
\]
Set
\[
 c=\frac{p+q}{2},\quad u=\frac{p-q}{2},\quad a=\norm{u},\quad y=x-c.
\]
Then
\[
 F_{p,q}(x)=(\norm{y}^2+a^2)^2-4\ip{y}{u}^2.
\]
Looking at this form, one can guess why the two-centre case is explicitly computable:
apart from the radial term, only the direction \(u\) plays a role.  The
quantity controlling the Hessian-positive region will be
\[
D(y)=3\|y\|^4+4\langle y,u\rangle^2-2a^2\|y\|^2-a^4.
\]


\begin{lemma}\label{lem:two-centre-hessian}
One has
\[
 \nabla F_{p,q}(x)=4(\norm{y}^2+a^2)y-8\ip{y}{u}u
\]
and
\[
 HF_{p,q}(x)=4(\norm{y}^2+a^2)I+8(yy^T-uu^T).
\]
\end{lemma}

\begin{proof}
Since $p=c+u$ and $q=c-u$, we have
\[
 F_{p,q}(x)=\norm{y-u}^2\norm{y+u}^2
 =\bigl(\norm{y}^2+a^2-2\ip{y}{u}\bigr)
  \bigl(\norm{y}^2+a^2+2\ip{y}{u}\bigr),
\]
which gives the displayed normal form. Differentiating with respect to $y$
gives the gradient formula, and a second differentiation gives the Hessian
formula. The translation $x\mapsto y=x-c$ does not change these derivatives.
\end{proof}

The gradient formula already determines the critical set.  Any component of
\(y\) orthogonal to \(u\) is multiplied by the positive factor
\(\|y\|^2+a^2\), so a critical point must lie on the line through the two
centres.

\begin{proposition}\label{prop:two-centre-critical-set}(see also \cite[Example 2.1]{BrojbeanuPintea-JCA})
The critical set of $F_{p,q}$ is
\[
 C(F_{p,q})=\{p,c,q\},
\]
and the corresponding critical values are $0,a^4,0$.
\end{proposition}

\begin{proof}
If $\nabla F_{p,q}(x)=0$, then
\[
 (\norm{y}^2+a^2)y=2\ip{y}{u}u.
\]
The component of $y$ orthogonal to $u$ must vanish, since
$\norm{y}^2+a^2>0$. Hence $y=tu$ for some $t\in\R$. Substitution gives
\[
 (t^2+1)t=2t,
\]
so $t\in\{-1,0,1\}$. Thus $x=c+tu$ is one of $p,c,q$, and direct evaluation
gives the stated critical values.
\end{proof}

For the computation of \(h_{\max}\), the critical points alone are not
enough; we need the full Hessian-positive region. The next result is a generalization to arbitrary dimension of \cite[Remark 3.2]{BrojbeanuPintea-JCA}. 


\begin{theorem}\label{thm:two-centre-hess-region}
For $n\ge2$,
\[
 HF_{p,q}(x)\succ0
 \quad\Longleftrightarrow\quad
 3\norm{y}^4+4\ip{y}{u}^2-2a^2\norm{y}^2>a^4.
\]
\end{theorem}

\begin{proof}
Let $e_1=u/a$. If $y$ is not parallel to $u$, choose a unit vector $e_2$
so that $y\in\operatorname{span}\{e_1,e_2\}$. If $y$ is parallel to $u$,
choose any unit vector $e_2\perp e_1$. Write $y=se_1+te_2$, with $t=0$
in the parallel case.

On $\operatorname{span}\{e_1,e_2\}^{\perp}$, the two rank-one terms in the
Hessian vanish, so $HF_{p,q}$ acts by the positive scalar
$4(\norm{y}^2+a^2)$. On $\operatorname{span}\{e_1,e_2\}$, the Hessian block is
\[
 4\begin{pmatrix}
 3s^2+t^2-a^2 & 2st\\
 2st & s^2+3t^2+a^2
 \end{pmatrix}.
\]
The determinant of the inner matrix is
\[
 3(s^2+t^2)^2+4a^2s^2-2a^2(s^2+t^2)-a^4
 =D(y),
\]
because $\ip{y}{u}=as$. The trace of the full $2\times2$ block is
$16(s^2+t^2)=16\norm{y}^2$. Thus, if $D(y)>0$, then $y\ne0$ and this trace is
positive; for a real symmetric $2\times2$ matrix this is equivalent to
positive definiteness of the block. Conversely, a positive definite block has
positive determinant, hence $D(y)>0$. Combining this with the positive
orthogonal directions proves the criterion. The same block decomposition gives
the useful determinant identity
\[
 \det HF_{p,q}(x)=4^n(\norm{y}^2+a^2)^{n-2}D(y),
\]
but the positive-definiteness criterion above uses the block decomposition, not
determinant positivity alone in higher dimension.
\end{proof}


\begin{thebibliography}{99}
\bibitem[]{BrojbeanuPintea-JCA} Brojbeanu, A.; Pintea, C.. Properties of the Level Sets of Some Products of Functions. Journal of Convex Analysis 30 (2023) 271--294
\end{thebibliography}
\end{document}
